\chapter{Methods and Software Implementation}
\section{Entry-aware splitting}
For a candidate split at a node, let $t_1,\ldots,t_K$ be its distinct event
times. Let $d_k$ and $d_{1k}$ denote event counts in the parent and left
daughter, and let $Y_k$ and $Y_{1k}$ denote their risk-set sizes. The
delayed-entry log-rank numerator and variance term are
\begin{align}
 U&=\sum_{k:Y_k>0}\left(d_{1k}-d_k\frac{Y_{1k}}{Y_k}\right),\\
 V&=\sum_{k:Y_k>1}
 \frac{Y_{1k}(Y_k-Y_{1k})d_k(Y_k-d_k)}{Y_k^2(Y_k-1)}.
\end{align}
Candidate splits are compared using the magnitude of the standardized
statistic when its variance is positive. Entry awareness changes the
risk-set counts through Equation~\eqref{eq:risk}; observed event counts are
unchanged. These are the risk-set modifications specified in the
project outline \citep{projectoutline}.

\section{Terminal-node estimates and returned predictions}
For terminal node $h$, let $d_h(t_k)$ be its event count and $Y_h(t_k)$ its
entry-aware risk-set size. The node estimates are
\begin{align}
 \widehat H_h(t)&=\sum_{t_k\leq t}\frac{d_h(t_k)}{Y_h(t_k)},\\
 \widehat S_h(t)&=\prod_{t_k\leq t}
       \left(1-\frac{d_h(t_k)}{Y_h(t_k)}\right).
\end{align}
Only node event times with a nonzero risk set contribute. A new subject is
routed to node $h_b(x)$ in tree $b$ using its covariates. The package returns
averages of the tree-specific cumulative hazards and product-limit curves:
\begin{align}
 \widehat H_e(t\mid x)&=\frac{1}{B}\sum_{b=1}^{B}\widehat H_{h_b(x)}(t),\\
 \widehat S_e(t\mid x)&=\frac{1}{B}\sum_{b=1}^{B}\widehat S_{h_b(x)}(t).
\end{align}
The subscript $e$ denotes the ensemble. For a training subject,
out-of-bag (OOB) predictions use only trees whose bootstrap samples
excluded that subject, averaging over the number of available OOB trees.

For RSF, the experiments use the package's returned survival predictions,
$\widehat S_e(t\mid x)$, obtained by averaging the tree-specific
Kaplan--Meier curves.

Entry times determine the risk sets used during tree construction and
terminal-node estimation. To obtain a new subject's survival curve,
the fitted trees use the predictor values $x$ to select terminal nodes.
The subject's entry time is not required for this step. The returned
curve describes survival from the original time origin. Entry time is
used when conditioning the prediction on survival to entry or evaluating
it with an entry-aware outcome-based score.

\section{Integration, validation, and supported scope}
The user supplies one entry time per original input row through
\texttt{rfsrc(..., entry.time = L)}. The interface aligns entries with rows
retained after missing-data processing and checks that retained entries are
finite, nonnegative, and no later than follow-up. Positive entry with a
missing follow-up time is rejected; missing predictors remain imputable.
The C implementation uses a lower-bound lookup to map entries onto the
master time grid, without requiring entries to coincide with exit times.

Entry times are stored on the grow object and saved forest, and are passed
explicitly during restoration and prediction. Refitting on subsets or
bootstrap samples subsets entries with the observations. Tuning aligns
training and holdout entries and uses holdout subjects' entries for the LT
concordance comparison. Training outcomes used to estimate nuisance
distributions are kept separate from evaluation outcomes.

Positive entries are supported with the single-event default log-rank rule
and with outcome-independent random splitting. Other rules, including
log-rank-score, Brier-gradient, and custom survival rules, are rejected with
positive entries because their split statistics have not been fully adapted.
Allowing random splits does not establish general LT validity for every
outcome family. The thesis evaluates single-event log-rank forests.

For nonnegative follow-up, setting all entries to zero reduces
Equation~\eqref{eq:risk} to the ordinary risk indicator. Ordinary calls can
then use the existing fast concordance algorithms. Positive-entry
concordance currently uses pairwise enumeration. Zero-entry agreement is a
necessary compatibility check, not a proof of correctness for positive entry.

\section{Evaluation targets}

The experiments evaluate both survival-curve accuracy and risk ranking.
MAD compares predicted survival with the known source-population curve
at independently sampled test covariates. Concordance, Brier score,
and AUC instead use the observed training outcomes: RSF contributes
OOB predictions, while Cox contributes in-sample predictions.
These outcome-based comparisons therefore use different validation
designs across model families.

Write $\widehat S_i(t)=\widehat S(t\mid X_i)$ for the prediction used
for subject $i$. Both fitting phases are scored using the same actual
entry times, even when entry is ignored during fitting.

\subsection{Concordance}

Concordance measures whether higher predicted risk agrees with earlier
observed failure. We summarize each curve over an evaluation grid
$\tau_1,\ldots,\tau_M$ using
\[
 r_i=\sum_{g=1}^{M}\{1-\widehat S_i(\tau_g)\}.
\]
Larger values indicate higher predicted risk. For an observed failure
with $\delta_i=1$ and $U_i<U_j$, a pair is comparable only if
$L_j\leq U_i$: subject $j$ must already have entered by subject $i$'s
failure time. Comparable pairs receive equal weight, with ties handled
by the package's existing conventions. We report the error $1-C$,
so lower values indicate better ranking.

\subsection{Entry-conditioned Brier score}

The Brier score measures squared error in predicted survival
probabilities. For subjects entered by time $t$, the implementation uses
\[
 q_i(t)=\frac{\widehat S_i(t)}{\widehat S_i(L_i^-)},
 \qquad t\geq L_i.
\]
Here $\widehat S_i(L_i^-)$ is predicted survival immediately before
entry and must be positive.
\alignmentnote{Section 9.2 uses $\widehat S_i(t)$. Should observation
conditioning use the ratio $q_i(t)$ above, or retain $\widehat S_i(t)$
as written in the outline?}

For background on censoring-adjusted Brier scoring, see
\citet[Section~4.1]{fuSimonoffSupplement}; related dynamic Brier methods
are documented in \citet{jmbayes2brier}. The entry-specific construction
below is our interpretation of Section 9.2 \citep{projectoutline}.

Let $\mathcal E_t$ contain subjects with $L_i\leq t$ and a finite
$q_i(t)$, and let $n_t=|\mathcal E_t|$. Write $\widehat G(t)$ for the
pooled censoring-survival estimate obtained using delayed-entry risk
sets. The implemented score is
\begin{align}
 \widehat{\mathrm{BS}}_{\mathrm{entry}}(t)
 =\frac{1}{n_t}\sum_{i\in\mathcal E_t}
 \Bigg[&
 \delta_i\ind(U_i\leq t)q_i(t)^2
 \frac{\widehat G(L_i^-)}{\widehat G(U_i^-)}
 \nonumber\\
 &+\ind(U_i>t)\{1-q_i(t)\}^2
 \frac{\widehat G(L_i^-)}{\widehat G(t)}
 \Bigg].
\end{align}

The first term scores observed failures by $t$; the second scores
subjects followed beyond $t$ without failure. The ratios are inverse
censoring-probability weights conditional on reaching entry.
Subjects censored by $t$ contribute zero but remain in $n_t$;
subjects not yet entered are excluded.

The pooled model assumes a common censoring hazard and requires
independent censoring and positive observation probabilities over
the scoring horizon. Some simulated censoring mechanisms depend
on entry or covariates, so the resulting comparisons are interpreted
descriptively under this working model.

We summarize the Brier curve by its unnormalized time integral,
reported as CRPS:
\[
 \mathrm{CRPS}=\int_{\tau_1}^{\tau_M}
       \widehat{\mathrm{BS}}_{\mathrm{entry}}(t)\,dt.
\]
The code approximates this integral by the trapezoidal rule.
Its scale depends on the evaluation horizon. If a Brier grid point
is undefined, the integrated score is also undefined.

\subsection{Additional time-dependent AUC}

At time $t$, AUC compares subjects with observed failures by $t$
(cases) against subjects followed beyond $t$ (controls). Using the
entered-subject set $\mathcal E_t$ defined above, let
\[
 \mathcal D_t=\{i\in\mathcal E_t:\delta_i=1,\ U_i\leq t\},
 \qquad
 \mathcal A_t=\{j\in\mathcal E_t:U_j>t\}.
\]
Define the time-specific risk marker $a_i(t)=1-q_i(t)$ and weights
\[
 w_i^{\mathrm{case}}
 =\frac{\widehat G(L_i^-)}{\widehat G(U_i^-)},
 \qquad
 w_j^{\mathrm{control}}(t)
 =\frac{\widehat G(L_j^-)}{\widehat G(t)}.
\]
The implemented AUC is
\begin{align}
 \widehat{\mathrm{AUC}}_{\mathrm{entry}}(t)
 &=
 \frac{
   \displaystyle\sum_{i\in\mathcal D_t}
   \sum_{j\in\mathcal A_t}
   w_i^{\mathrm{case}}w_j^{\mathrm{control}}(t)
   \left[
    \ind\{a_i(t)>a_j(t)\}
    +\tfrac12\ind\{a_i(t)=a_j(t)\}
   \right]
 }{
   \displaystyle
   \left(\sum_{i\in\mathcal D_t}w_i^{\mathrm{case}}\right)
   \left(\sum_{j\in\mathcal A_t}w_j^{\mathrm{control}}(t)\right)
 }.
\end{align}
The denominator is the total weight of all case--control pairs.
The numerator counts correctly ranked pairs with those weights,
giving half credit for tied markers. AUC is undefined when either
group is empty.

We report $1-\overline{\mathrm{AUC}}$, where the average is calculated
by trapezoidal integration over defined grid points and divided by
their time span. This descriptive average can bridge undefined interior
points. The censoring assumptions are the same as for Brier scoring.
\alignmentnote{Should this additional AUC definition be included
in the evaluation? The project outline does not specify an AUC estimator.}
